% ============ 第6章 §6.1 Littlewood--Paley 理论 (原书 pp.123-130) ============

这里, 作为奇异积分算子有界性理论的应用, 说明 $L^{p}$ 范数的有用改写是可能的. 这一应用称为 Littlewood--Paley 理论, 由此可以考察种种算子的性质.

\section{Littlewood--Paley 理论}\label{sec:6.1}

对 $\varphi \in \mathcal{S}(\mathbb{R}^{n})$ 与 $f \in \mathcal{S}'(\mathbb{R}^{n})$, 令 $\varphi(D)f = \mathcal{F}^{-1}[\varphi\mathcal{F}f] \in C^{\infty}(\mathbb{R}^{n})$. 这与定义为 $\varphi(D)f = (2\pi)^{-\frac{n}{2}}\mathcal{F}^{-1}\varphi * f$ 是同一回事. 另外, 对 $\varphi \in \mathcal{S}(\mathbb{R}^{n})$ 与 $j \in \mathbb{Z}$, 令 $\varphi_{j} = \varphi(2^{-j}\,\cdot\,)$.

\begin{theorem}\label{thm:6.1}
设 $1 < p < \infty$. 设只依赖于 $|x|$ 的实值函数 $\varphi \in \mathcal{S}(\mathbb{R}^{n})$ 满足 $\supp(\varphi) \subset B(4) \setminus B(1)$ 与 $\sum_{j=-\infty}^{\infty}\varphi_{j}^{2} = \chi_{\mathbb{R}^{n}\setminus\{0\}}$. 此时, 若 $f \in L^{p}(\mathbb{R}^{n})$, 则
\[
\|f\|_{L^{p}} \sim \left\|\left(\sum_{j=-\infty}^{\infty}|\varphi_{j}(D)f|^{2}\right)^{1/2}\right\|_{L^{p}}
\]
成立.
\end{theorem}

\begin{note}\label{note:6.1}
考察 $\sum_{j=-\infty}^{\infty}\varphi_{j}^{2} = \chi_{\mathbb{R}^{n}\setminus\{0\}}$ 这个条件.

(1) 由集合特征函数的定义可知, 这个条件在 $x \neq 0$ 时, 可读作
\[
\sum_{j=-\infty}^{\infty}\varphi_{j}(x)^{2} = 1,
\]
在 $x = 0$ 时, 可读作
\[
\sum_{j=-\infty}^{\infty}\varphi_{j}(x)^{2} = 0.
\]

(2) 因为 $\varphi(0) = 0$, 所以 $\varphi_{j}(0) = 0$. 因此, (1) 的 $x = 0$ 情形的条件是平凡的.

(3) 想考察是否真的存在这样的函数. 如从在 $t > 0$ 时, 用次数低于 $k$ 次的多项式 $p_{k}$ 可以表示为
\begin{equation}\tag{6.1}
\varphi^{(k)}(t) = \frac{p_{k}(t)}{t^{2k}}e^{-1/t}
\label{eq:6.1}
\end{equation}
的形式可以看出, 由 $\lim_{t \to \infty}t^{k}e^{-t} = 0$ ($k = 0, 1, 2, \ldots$), 对 $t \leq 0$ 满足 $\varphi(t) = 0$, 对 $t > 0$ 满足 $\varphi(t) = e^{-1/t}$ 的函数 $\varphi$ 是 $C^{\infty}$ 级函数. 也就是说, 无论微分多少次都可以. 因此,
\[
\Phi(x) = \frac{\varphi(1 - |x|^{2})}{\varphi(1 - |x|^{2}) + \varphi(|x|^{2})}
\]
是在 $|x| < 1$ 取正值, 在 $|x| \geq 1$ 取 $0$ 的, 定义在全空间 $\mathbb{R}^{n}$ 上的 $C_{c}^{\infty}$ 级函数. 另外,
\[
\Psi(x) = \frac{\varphi(9|x|^{2} - 1)}{\varphi(1 - |x|^{2}) + \varphi(9|x|^{2} - 1)}
\]
是在 $|x| > \frac{1}{3}$ 取正值, 在 $|x| \leq \frac{1}{3}$ 取 $0$ 的, 定义在全空间 $\mathbb{R}^{n}$ 上的 $C^{\infty}$ 级函数. 因此,
\[
\Theta(x) = \Phi(x)\Psi(x) \quad (x \in \mathbb{R}^{n})
\]
是仅在 $\frac{1}{3} < |x| < 1$ 取正值的, 定义在全空间 $\mathbb{R}^{n}$ 上的 $C^{\infty}$ 级函数. 于是, 对 $x \in \mathbb{R}^{n}$, 令
\[
\zeta(x) = \sqrt{\sum_{j=-\infty}^{\infty}\Theta(2^{-j}x)^{2}},
\]
若 $x \neq 0$, 则
\[
\varphi(x) = \frac{1}{\zeta(x)}\Theta\left(\frac{2}{7}x\right)
\]
通过补充定义 $\varphi(0) = 0$, 成为 $C^{\infty}$ 级函数, 支集包含在 $B(4) \setminus B(1)$ 中, 且满足 $\sum_{j=-\infty}^{\infty}\varphi_{j}^{2} = \chi_{\mathbb{R}^{n}\setminus\{0\}}$.
\end{note}

$p = 2$ 的情形可以由 Fourier 变换的性质得到结论.

\begin{lemma}\label{lem:6.1}
若 $p = 2$, 则定理~\ref{thm:6.1} 成立. 更详细地, 对一切 $f \in L^{2}(\mathbb{R}^{n})$, 有
\begin{equation}\tag{6.2}
\left\|\left(\sum_{j=-\infty}^{\infty}|\varphi_{j}(D)f|^{2}\right)^{1/2}\right\|_{L^{2}} = \|f\|_{L^{2}}
\label{eq:6.2}
\end{equation}
成立.
\end{lemma}

\begin{proof}
由单调收敛定理,
\[
\left\|\left(\sum_{j=-\infty}^{\infty}|\varphi_{j}(D)f|^{2}\right)^{1/2}\right\|_{L^{2}} = \left(\sum_{j=-\infty}^{\infty}(\|\varphi_{j}(D)f\|_{L^{2}})^{2}\right)^{1/2}.
\]
由于 $\mathcal{F}^{-1}$ 由 Plancherel 定理保持 $L^{2}$ 范数,
\[
\left\|\left(\sum_{j=-\infty}^{\infty}|\varphi_{j}(D)f|^{2}\right)^{1/2}\right\|_{L^{2}} = \left(\sum_{j=-\infty}^{\infty}(\|\varphi_{j}\mathcal{F}f\|_{L^{2}})^{2}\right)^{1/2}.
\]
再次由单调收敛定理,
\[
\left\|\left(\sum_{j=-\infty}^{\infty}|\varphi_{j}(D)f|^{2}\right)^{1/2}\right\|_{L^{2}} = \left\|\left(\sum_{j=-\infty}^{\infty}|\varphi_{j}\mathcal{F}f|^{2}\right)^{1/2}\right\|_{L^{2}}.
\]
这里使用 $\sum_{j=-\infty}^{\infty}\varphi_{j}^{2} = \chi_{\mathbb{R}^{n}\setminus\{0\}}$, 得
\[
\left\|\left(\sum_{j=-\infty}^{\infty}|\varphi_{j}(D)f|^{2}\right)^{1/2}\right\|_{L^{2}} = \|\mathcal{F}f\|_{L^{2}}.
\]
因此, 再次由 Plancherel 定理, 由于 $\mathcal{F}$ 保持 $L^{2}$ 范数, 得到式~\eqref{eq:6.2}.
\end{proof}

证明 $p \neq 2$ 的情形时, 下面的引理很管用.

\begin{lemma}\label{lem:6.2}
对 $a > 0$, 由
\[
\Xi(a) \equiv \sum_{j=-\infty}^{\infty}(2^{j}a)^{n}(1 + 2^{j}a)^{-n-1}
\]
给出的函数 $\Xi: (0,\infty) \to \mathbb{R}$ 是有界的.
\end{lemma}

\begin{proof}
假设 $1 < a \leq 2$, 则
\begin{align*}
\Xi(a) &\leq 2^{n}\sum_{j=-\infty}^{\infty}2^{jn}(1 + 2^{j})^{-n-1} \\
&\leq 2^{n}\sum_{j=0}^{\infty}2^{jn}\cdot(2^{j})^{-n-1} + 2^{n}\sum_{j=-\infty}^{-1}2^{jn} < \infty.
\end{align*}
由于 $\Xi(2a) = \Xi(a)$, $\Xi(0) = 0$, 所以对一切 $a \geq 0$, $\Xi(a) \lesssim 1$ 成立.
\end{proof}

\begin{proof}
来证明 $p \neq 2$ 的情形下的定理~\ref{thm:6.1}. 首先证明式~\eqref{eq:6.3}:
\begin{equation}\tag{6.3}
\left\|\left(\sum_{j=-\infty}^{\infty}|\varphi_{j}(D)f|^{2}\right)^{1/2}\right\|_{L^{p}} \lesssim \|f\|_{L^{p}}
\label{eq:6.3}
\end{equation}
关于 $j$ 适当地截断, 之后用 Fatou 引理令 $J \to \infty$, 从而只要证明与自然数 $J$ 无关的一致估计
\[
\left\|\left(\sum_{j=-J}^{J}|\varphi_{j}(D)f|^{2}\right)^{1/2}\right\|_{L^{p}} \lesssim \|f\|_{L^{p}}
\]
即可.

设 $\{r_{j}\}_{j=1}^{\infty}$ 为先前考虑过的 Rademacher 序列. 将 $\{r_{j}\}_{j=1}^{\infty}$ 不重复地重新排列, 使得 $\{\tilde{r}_{j}\}_{j=-\infty}^{\infty} = \{r_{j}\}_{j=1}^{\infty}$. 由此,
\[
\left\|\left(\sum_{j=-J}^{J}|\varphi_{j}(D)f|^{2}\right)^{1/2}\right\|_{L^{p}} \lesssim \left\|\sum_{j=-J}^{J}\tilde{r}_{j} \otimes \varphi_{j}(D)f\right\|_{L^{p}([0,1]_{t} \times \mathbb{R}^{n}_{x})}.
\]
所以, 只要证明固定 $t \in [0,1]$ 时, 算子
\[
T_{t}: f \mapsto \sum_{j=-J}^{J}\tilde{r}_{j}(t)\mathcal{F}^{-1}\varphi_{j} * f
\]
是 $L^{p}$ 有界的, 且其算子范数可用与 $t$ 无关的数来估计即可. 对应的积分核 $K_{t}$ 具体呈如下形状.
\[
K_{t}(z) = \sum_{j=-J}^{J}\tilde{r}_{j}(t)\mathcal{F}^{-1}\varphi_{j}(z).
\]

(1) 关于 $T_{t}$ 的 $B(L^{2}(\mathbb{R}^{n}))$ 范数 $\|T_{t}\|_{L^{2} \to L^{2}}$, 由
\[
A \equiv \sup_{\xi \in \mathbb{R}^{n}}\sum_{j=-\infty}^{\infty}|\varphi_{j}(\xi)| \lesssim 1
\]
与 Plancherel 定理, 可得 $\|T_{t}\|_{L^{2} \to L^{2}} \leq A$.

(2) 因为 $\mathcal{F}^{-1}\varphi_{j} = 2^{jn}\mathcal{F}\varphi(2^{j}\,\cdot\,)$, 所以使用引理~\ref{lem:6.2} 的函数 $\Xi$,
\[
\begin{aligned}
|x-y|^{n}|K_{t}(z)| &\leq \sum_{j=-\infty}^{\infty}|x-y|^{n}|\mathcal{F}^{-1}\varphi_{j}(x-y)| \\
&\lesssim \sum_{j=-\infty}^{\infty}(2^{j}|x-y|)^{n}(1 + 2^{j}|x-y|)^{-n-1} \\
&= \Xi(|x-y|).
\end{aligned}
\]
因此, 由引理~\ref{lem:6.2}, 可得
\[
|x-y|^{n}|K_{t}(z)| \lesssim 1.
\]

(3) 因为 $\nabla K_{t}(z) = 2^{j(n+1)}(\nabla\mathcal{F}^{-1}\varphi)(2^{j}\,\cdot\,)$, 所以使用 $|\nabla\mathcal{F}^{-1}\varphi(z)| \lesssim (1+|z|)^{-n-2}$, 与刚才一样, 可得
\[
\begin{aligned}
|x-y|^{n+1}|\nabla K_{t}(z)| &\leq \sum_{j=-\infty}^{\infty}|x-y|^{n+1}|\nabla(\mathcal{F}^{-1}\varphi_{j})(x-y)| \\
&\lesssim \sum_{j=-\infty}^{\infty}(2^{j}|x-y|)^{n+1}(1 + 2^{j}|x-y|)^{-n-2} \\
&\lesssim 1.
\end{aligned}
\]

由以上可知, $T_{t}$ 是一致地 $L^{p}$ 有界的.

来证明反向的不等式
\[
\left\|\left(\sum_{j=-\infty}^{\infty}|\varphi_{j}(D)f|^{2}\right)^{1/2}\right\|_{L^{p}} \gtrsim \|f\|_{L^{p}}.
\]
暂时假设 $f \in L^{2}(\mathbb{R}^{n}) \cap L^{p}(\mathbb{R}^{n})$. 由对偶性,
\[
\|f\|_{L^{p}} = \sup\left\{\left|\int_{\mathbb{R}^{n}}f(x)g(x)\dd x\right|: g \in L^{p'}(\mathbb{R}^{n}), \|g\|_{L^{p'}} \leq 1\right\}.
\]
因为 $\mathcal{S}(\mathbb{R}^{n})$ 在 $L^{p}(\mathbb{R}^{n})$ 中稠密,
\[
\|f\|_{L^{p}} = \sup\left\{\left|\int_{\mathbb{R}^{n}}f(x)g(x)\dd x\right|: g \in \mathcal{S}(\mathbb{R}^{n}), \|g\|_{L^{p'}} \leq 1\right\}.
\]
基于这一点, 固定 $g \in \mathcal{S}(\mathbb{R}^{n})$. 因为 $f, g \in L^{2}(\mathbb{R}^{n})$, 由引理~\ref{lem:6.1},
\[
\begin{aligned}
\int_{\mathbb{R}^{n}}\sum_{j=-\infty}^{\infty}|\varphi_{j}(D)f(x)\varphi_{j}(D)g(x)|\dd x &\leq \sqrt{\int_{\mathbb{R}^{n}}\sum_{j=-\infty}^{\infty}|\varphi_{j}(D)f(x)|^{2}\dd x\int_{\mathbb{R}^{n}}\sum_{j=-\infty}^{\infty}|\varphi_{j}(D)g(x)|^{2}\dd x} \\
&\leq \|f\|_{L^{2}}\|g\|_{L^{2}} < \infty.
\end{aligned}
\]
因此, 由
\[
\sum_{j=-\infty}^{\infty}\varphi_{j}^{2} = \chi_{\mathbb{R}^{n}\setminus\{0\}},
\]
可变形为
\[
\begin{aligned}
\int_{\mathbb{R}^{n}}f(x)\overline{g(x)}\dd x &= \int_{\mathbb{R}^{n}}\mathcal{F}f(\xi)\overline{\mathcal{F}g(\xi)}\dd\xi \\
&= \sum_{j=-\infty}^{\infty}\int_{\mathbb{R}^{n}}\varphi_{j}(\xi)\mathcal{F}f(\xi)\overline{\varphi_{j}(\xi)\mathcal{F}g(\xi)}\dd\xi \\
&= \sum_{j=-\infty}^{\infty}\int_{\mathbb{R}^{n}}\varphi_{j}(D)f(x)\overline{\varphi_{j}(D)g(x)}\dd x \\
&= \int_{\mathbb{R}^{n}}\left(\sum_{j=-\infty}^{\infty}\varphi_{j}(D)f(x)\overline{\varphi_{j}(D)g(x)}\right)\dd x.
\end{aligned}
\]
这里, 由 Cauchy--Schwarz 不等式,
\[
\begin{aligned}
\left|\int_{\mathbb{R}^{n}}f(x)\overline{g(x)}\dd x\right| &\leq \int_{\mathbb{R}^{n}}\sum_{j=-\infty}^{\infty}|\varphi_{j}(D)f(x)\varphi_{j}(D)g(x)|\dd x \\
&\leq \int_{\mathbb{R}^{n}}\left(\sum_{j=-\infty}^{\infty}|\varphi_{j}(D)f(x)|^{2}\right)^{1/2}\left(\sum_{j=-\infty}^{\infty}|\varphi_{j}(D)g(x)|^{2}\right)^{1/2}\dd x.
\end{aligned}
\]
进而, 通过使用 Hölder 不等式, 可得
\[
\begin{aligned}
\left|\int_{\mathbb{R}^{n}}f(x)\overline{g(x)}\dd x\right| &\leq \left\|\left(\sum_{j=-\infty}^{\infty}|\varphi_{j}(D)f|^{2}\right)^{1/2}\right\|_{L^{p}}\left\|\left(\sum_{j=-\infty}^{\infty}|\varphi_{j}(D)g|^{2}\right)^{1/2}\right\|_{L^{p'}}.
\end{aligned}
\]
这里, 因为 $g \in L^{p'}(\mathbb{R}^{n})$,
\[
\left\|\left(\sum_{j=-\infty}^{\infty}|\varphi_{j}(D)g|^{2}\right)^{1/2}\right\|_{L^{p'}} \lesssim \|g\|_{L^{p'}} \leq 1
\]
可得. 因此, 若 $f \in L^{2}(\mathbb{R}^{n}) \cap L^{p}(\mathbb{R}^{n})$, 则
\begin{equation}\tag{6.4}
\left\|\left(\sum_{j=-\infty}^{\infty}|\varphi_{j}(D)f|^{2}\right)^{1/2}\right\|_{L^{p}} \gtrsim \|f\|_{L^{p}}
\label{eq:6.4}
\end{equation}
就得到证明.

最后, 考虑把假设 $f \in L^{2}(\mathbb{R}^{n}) \cap L^{p}(\mathbb{R}^{n})$ 放宽为 $f \in L^{p}(\mathbb{R}^{n})$. 设 $f \in L^{p}(\mathbb{R}^{n})$. 因为 $L^{2}(\mathbb{R}^{n}) \cap L^{p}(\mathbb{R}^{n})$ 在 $L^{p}(\mathbb{R}^{n})$ 中稠密, 可以取列 $f_{k} \in L^{2}(\mathbb{R}^{n}) \cap L^{p}(\mathbb{R}^{n})$, $k = 1, 2, \ldots$, 使得 $\lim_{k \to \infty}f_{k} = f$ 在 $L^{p}(\mathbb{R}^{n})$ 中成立. 由三角不等式与式~\eqref{eq:6.4},
\[
\begin{aligned}
\|f\|_{L^{p}} &\leq \|f - f_{k}\|_{L^{p}} + \|f_{k}\|_{L^{p}} \\
&\lesssim \|f - f_{k}\|_{L^{p}} + \left\|\left(\sum_{j=-\infty}^{\infty}|\varphi_{j}(D)f_{k}|^{2}\right)^{1/2}\right\|_{L^{p}}.
\end{aligned}
\]
进而, 由三角不等式与式~\eqref{eq:6.3},
\[
\begin{aligned}
\|f\|_{L^{p}} &\leq \|f - f_{k}\|_{L^{p}} + \left\|\left(\sum_{j=-\infty}^{\infty}|\varphi_{j}(D)[f_{k} - f]|^{2}\right)^{1/2}\right\|_{L^{p}} \\
&\quad + \left\|\left(\sum_{j=-\infty}^{\infty}|\varphi_{j}(D)f|^{2}\right)^{1/2}\right\|_{L^{p}} \\
&\lesssim \|f - f_{k}\|_{L^{p}} + \left\|\left(\sum_{j=-\infty}^{\infty}|\varphi_{j}(D)f|^{2}\right)^{1/2}\right\|_{L^{p}},
\end{aligned}
\]
所以令 $k \to \infty$, 就得到关于 $f$ 的结论式
\[
\left\|\left(\sum_{j=-\infty}^{\infty}|\varphi_{j}(D)f|^{2}\right)^{1/2}\right\|_{L^{p}} \gtrsim \|f\|_{L^{p}}.
\]
\end{proof}

作为 Littlewood--Paley 理论的简单应用, 给出下面这一结果. 用
\[
\mathcal{S}_{\infty}(\mathbb{R}^{n}) = \bigcap_{\alpha \in \mathbb{N}_{0}^{n}}\left\{f \in \mathcal{S}(\mathbb{R}^{n}): \int_{\mathbb{R}^{n}}x^{\alpha}f(x)\dd x = 0\right\}
\]
定义 $\mathcal{S}_{\infty}(\mathbb{R}^{n})$. $\mathcal{S}_{\infty}(\mathbb{R}^{n})$ 的元素有时也称为 Lizorkin 分布.

\begin{corollary}\label{cor:6.1}
若 $1 < p < \infty$, 则 $\mathcal{S}_{\infty}(\mathbb{R}^{n})$ 在 $L^{p}(\mathbb{R}^{n})$ 中稠密.
\end{corollary}

\begin{proof}
已知 $C_{c}^{\infty}(\mathbb{R}^{n})$ 在 $L^{p}(\mathbb{R}^{n})$ 中稠密, 所以只要证明任意 $f \in C_{c}^{\infty}(\mathbb{R}^{n})$ 可以用 $\mathcal{S}_{\infty}(\mathbb{R}^{n})$ 的拓扑来近似即可. 设 $\psi \in C^{\infty}(\mathbb{R}^{n})$ 满足 $\chi_{B(4)} \leq \psi \leq \chi_{B(8)}$.

由注意~\ref{note:3.2}, Hardy--Littlewood 极大算子的 $L^{p}$ 有界性以及 Lebesgue 收敛定理, 在 $L^{p}(\mathbb{R}^{n})$ 的拓扑中,
\[
\lim_{j \to \infty}\psi_{j}(D)f = (2\pi)^{-\frac{n}{2}}\lim_{j \to \infty}2^{jn}\mathcal{F}^{-1}(2^{j}\,\cdot\,) * f = f.
\]
另外, 由 Littlewood--Paley 理论,
\[
\limsup_{k \to -\infty}\|\psi_{k}(D)f\|_{L^{p}} \lesssim \limsup_{k \to -\infty}\left\|\left(\sum_{j=-\infty}^{\infty}|\varphi_{j}(D)\psi_{k}(D)f|^{2}\right)^{1/2}\right\|_{L^{p}}.
\]
这里, $\varphi_{j}(D)\psi_{k}(D)f$ 在 $j > k+2$ 时为 $0$, 在 $j < k-2$ 时与 $\varphi_{j}(D)f$ 一致, 所以可以分解为
\[
\begin{aligned}
\limsup_{k \to -\infty}\|\psi_{k}(D)f\|_{L^{p}} &\lesssim \limsup_{k \to -\infty}\left\{\left\|\left(\sum_{j=-\infty}^{k-2}|\varphi_{j}(D)f|^{2}\right)^{1/2}\right\|_{L^{p}}\right. \\
&\qquad \left. + \sum_{j=k-2}^{k+2}\|\varphi_{j}(D)\psi_{k}(D)f\|_{L^{p}}\right\}.
\end{aligned}
\]
这里, 由 Young 不等式,
\[
\begin{aligned}
\limsup_{k \to -\infty}\|\psi_{k}(D)f\|_{L^{p}} &\lesssim \limsup_{k \to -\infty}\left\{\left\|\left(\sum_{j=-\infty}^{k-2}|\varphi_{j}(D)f|^{2}\right)^{1/2}\right\|_{L^{p}}\right. \\
&\qquad \left. + \sum_{j=k-2}^{k+2}\|\varphi_{j}(D)f\|_{L^{p}}\right\},
\end{aligned}
\]
所以由积分的单调性, 可归纳为
\[
\limsup_{k \to -\infty}\|\psi_{k}(D)f\|_{L^{p}} \lesssim \limsup_{k \to -\infty}\left\|\left(\sum_{j=-\infty}^{k+2}|\varphi_{j}(D)f|^{2}\right)^{1/2}\right\|_{L^{p}}.
\]
因此,
\[
\left\|\left(\sum_{j=-\infty}^{k+2}|\varphi_{j}(D)f|^{2}\right)^{1/2}\right\|_{L^{p}} \lesssim \|f\|_{L^{p}} < \infty,
\]
所以由 Lebesgue 收敛定理, 当 $k \to -\infty$ 时, 得到 $\psi_{k}(D)f \to 0$. 这意味着
\[
f = \lim_{k \to \infty}(\psi_{k}(D)f - \psi_{-k}(D)f),
\]
且因为 $(\psi_{k}(D)f - \psi_{-k}(D)f) \in \mathcal{S}_{\infty}(\mathbb{R}^{n})$, 所以 $f$ 可以由 $\mathcal{S}_{\infty}(\mathbb{R}^{n})$ 的元素在 $L^{p}(\mathbb{R}^{n})$ 中近似.
\end{proof}

\begin{exercise}\label{exer:6.1}
对 $a > 0$, 证明由
\[
\Xi(a) \equiv \sum_{j=-\infty}^{\infty}\frac{1}{2^{-j}a^{-1} + 2^{j}a}
\]
给出的函数 $\Xi: (0,\infty) \to \mathbb{R}$ 是有界的.
\end{exercise}
